\documentclass[../main.tex]{subfiles}

\graphicspath{{../image/}{../figure/}}

\begin{document}

\chapter{补偿紧致性}\label{chapter:ch17-补偿紧致性}

双曲守恒律方程组的近似解可由多种方式生成：消失黏性方法、差分逼近、松弛格式等等。本章要讨论的主题是：解是否可构造为近似解序列的极限，而这些近似解序列只在某 $L^p$ 空间中有界。由于方程组是非线性的，困难在于：构造格式一般只在近似解序列强收敛时才相容，而所假设的 $L^p$ 界只能保证弱收敛——近似解可能发展出有限振幅的高频振荡，从而严重破坏相容性。我们的目标是证明：熵不等式可通过淬灭快速振荡而力挽狂澜，从而强制近似解强收敛。这一效应的某种迹象已在\cref{section:ch01-快速振荡}中提及。

本研究的主要工具将是 Young 测度的概念与补偿紧致性的泛函分析方法。前者自然引出极为一般的测度值解类，后者则被用于验证非线性把测度值解约化为传统解。由于该方法本质地依赖于熵耗散，它似乎主要适用于配备了丰富熵--熵流族的方程组，最显著的是标量守恒律以及恰由两个守恒律组成的方程组。尽管有此限制，该方法仍相当富有成效，不仅因为这类结构的重要方程组大量存在，还因为它对熵耗散的稳定作用以及非线性“冲突的”稳定--失稳行为提供了宝贵洞见。这些因素的不同表现已在早先各章中遇到过。

在该方法众多已知应用中，这里只呈现最简单者，涉及标量守恒律、两个守恒律的真非线性方程组，以及等熵弹性与气体动力学方程组。

\subfile{Chapter-17/section-01.tex}
\subfile{Chapter-17/section-02.tex}
\subfile{Chapter-17/section-03.tex}
\subfile{Chapter-17/section-04.tex}
\subfile{Chapter-17/section-05.tex}
\subfile{Chapter-17/section-06.tex}
\subfile{Chapter-17/section-07.tex}
\subfile{Chapter-17/section-08.tex}
\subfile{Chapter-17/section-09.tex}

\end{document}
